\section{Stochastic crowning}
\label{app:stochastic-crowning}

\subsection{Scope}
\label{app:stoch-scope}
A\(\Omega\)D stochastic accounting is count-normalized boundary accounting over a declared field, boundary, window, and fractal address.


\aodliteraturenote{Finite Markov kernels: Norris~\cite{norrisMarkov}, Levin--Peres--Wilmer~\cite{levinPeresWilmer}.}

\subsection{Setup}
\label{app:stoch-setup}
Declare a finite history family on the field and window:
\[
\mathcal H_{\mathcal F}(\omega).
\]
Each history has the form
\[
h=(h_{\mathrm{past}},h_{\mathrm{present}},h_{\mathrm{future}}).
\]

\subsection{Statement(s)}
\label{app:stoch-statements}

\subsubsection{History family}
\label{app:stoch-history-family}
The family \(\mathcal H_{\mathcal F}(\omega)\) is the declared finite support for stochastic crowning on the window.

\subsubsection{Crowning}
\label{app:stoch-crowning}
The crowning map selects the present:
\[
\operatorname{Crown}(h)=h_{\mathrm{present}}.
\]

\subsubsection{Alpha advantage}
\label{app:stoch-alpha-advantage}
At one declared observation restriction, partition histories into mutually
exclusive statuses: crowned-alpha and not yet returned, completed Omega
return, or unresolved. A history is counted in exactly one current status;
its earlier alpha event remains in provenance but is not counted again as a
second current history. With this partition, define
\[
N_\alpha=
\#\{h:\ h\text{ has current alpha status}\},
\]
\[
N_\Omega=
\#\{h:\ h\text{ has completed-return status}\},
\]
\[
N_0=
\#\{h:\ h\text{ remains unresolved}\}.
\]
The total count is
\[
N_{\mathrm{tot}}=N_\alpha+N_\Omega+N_0,
\qquad
N_{\mathrm{tot}}>0.
\]
The alpha advantage is
\begin{equation}
\Delta_\alpha(\omega)=
\frac{N_\alpha-N_\Omega}{N_{\mathrm{tot}}}.
\label{app:stochastic:eq:alpha-advantage}
\end{equation}

\subsubsection{Retained-history family}
\label{app:stoch-retained-history-family}
For channel \(e\), the retained-window indicator is
\[
\mathbf 1_{\mathrm{ret}}(e,\omega)=
\begin{cases}
1, & \mathrm{TTL}_e(\omega)\ge |\omega|_e,\\
0, & \mathrm{TTL}_e(\omega)< |\omega|_e.
\end{cases}
\]
The history retention indicator is
\[
\mathbf 1_{\mathrm{ret}}(h,\omega)=1
\]
when every retained channel required by \(h\) satisfies the declared TTL policy on \(\omega\). The same mutually exclusive status partition is used after filtering. The retained-history family is
\[
\mathcal H_{\mathcal F}^{\mathrm{ret}}(\omega)
=
\{h\in\mathcal H_{\mathcal F}(\omega):
\mathbf 1_{\mathrm{ret}}(h,\omega)=1\}.
\]
Define
\[
N_{\alpha}^{\mathrm{ret}}
=
\#\{h\in\mathcal H_{\mathcal F}^{\mathrm{ret}}(\omega):
h\text{ has current alpha status}\},
\]
\[
N_{\Omega}^{\mathrm{ret}}
=
\#\{h\in\mathcal H_{\mathcal F}^{\mathrm{ret}}(\omega):
h\text{ has completed-return status}\},
\]
\[
N_{\mathrm{tot}}^{\mathrm{ret}}
=
|\mathcal H_{\mathcal F}^{\mathrm{ret}}(\omega)|.
\]
The retained alpha advantage is defined when \(N_{\mathrm{tot}}^{\mathrm{ret}}>0\):
\begin{equation}
\Delta_{\alpha}^{\mathrm{ret}}(\omega)
=
\frac{N_{\alpha}^{\mathrm{ret}}-N_{\Omega}^{\mathrm{ret}}}{N_{\mathrm{tot}}^{\mathrm{ret}}}.
\label{app:stochastic:eq:ret-alpha-advantage}
\end{equation}

\subsection{Derivation}
\label{app:stoch-derivation}
For each history, let \(P_W(h)\) be the restriction of
\eqref{eq:temporal-sadar-burden} to its temporal cycle group.
The notation \(\mathbb E_{\mathcal H}\) denotes finite arithmetic average over the declared history family \(\mathcal H_{\mathcal F}(\omega)\):
\[
\mathbb E_{\mathcal H}[P_W]
=
\frac{1}{N_{\mathrm{tot}}}
\sum_{h\in\mathcal H_{\mathcal F}(\omega)}
P_W(h).
\]

\subsection{Falsification}
\label{app:stoch-falsification}
If a declared history family has \(N_{\mathrm{tot}}=0\), the alpha-advantage expression is undefined on that scope. If TTL filtering changes the retained support, retained counts must be recomputed before using \(\Delta_{\alpha}^{\mathrm{ret}}\).

\aodprovenance{\secref{subsec:crown},
\eqnref{app:stochastic:eq:alpha-advantage},
\eqnref{app:stochastic:eq:ret-alpha-advantage},
\secref{subsec:support-enclosure-retention-condition}.}

\aodremark{The probability terms are normalized counts on the declared history family.}
